Redes varejistas precisam escolher políticas de reposição para demandas heterogêneas. Esta dissertação propõe o \textit{Adaptive Inventory Policy Engine} (AIPE), que integra preparação, caracterização, calibração e seleção contextual. O \textit{Policy Selection Engine} (PSE) admite duas abordagens: associar o perfil informado a uma política ou estimar a política por série com aprendizado supervisionado. A avaliação usa sete candidatas, calibração no treino, escolha na validação de outras lojas e teste temporal reservado. Na Bahia, 145 séries em duas origens totalizam 4.350 ajustes; o PSE por perfil aumenta o custo médio em 1,93\% e o atendimento em 3,89 pontos percentuais, reduzindo de 119 para 100 os casos abaixo de 0,70 entre 283 com demanda positiva. No piloto M5, 30 séries das dez lojas em três origens totalizam 1.350 ajustes. O preditivo custa 13,9\% mais e atende 4,11 pontos menos que a global; incluir perfil não altera suas recomendações. Poucos exemplos por classe e instabilidade dos rótulos limitam o aprendizado. Um recuo exploratório ao perfil reduz a perda do preditivo, sem demonstrar superioridade sobre as referências. Oráculos retrospectivos e planejamento com informação perfeita separam erros de seleção, insuficiência do portfólio e restrições físicas iniciais. A contribuição é uma arquitetura e uma avaliação rastreável, com resultados favoráveis e limites explícitos. Coortes retrospectivas, calibrações específicas, janelas curtas e desenvolvimento no piloto exigem confirmação independente antes de atribuir vantagem geral ao AIPE.

\vspace{2em}
\textbf{Palavras-chave}: Gestão de estoques; Seleção de políticas; Demanda intermitente; Aprendizado supervisionado; Nível de serviço.
